\documentclass[10pt]{article}
\usepackage[utf8]{inputenc}
\usepackage[T1]{fontenc}
\usepackage{amsmath,amssymb,adjustbox}
\usepackage[margin=2cm]{geometry}
\begin{document}
\begin{center}\large Reducción y resolución formal\end{center}
\textbf{Modelo:} draft4\_case\_0\par
\[\adjustbox{max width=\linewidth}{$\displaystyle \text{Lagrangiano:}\qquad L=g{}^{a\,c}\,g{}^{b\,d}\,R{}_{a\,b\,c\,d} + \frac{2}{{\ell}^{2}}$}\]
\textbf{Ansatz:} draft4\_circular\_specialized\par
\[\adjustbox{max width=\linewidth}{$\displaystyle \text{Especialización de }\phi:\quad \phi=q\,\varphi$}\]
\textbf{Estado global:} hay familias verificadas y ramas pendientes\par
\[\adjustbox{max width=\linewidth}{$\displaystyle \text{Dominio de trabajo:}\qquad f\!\left(r\right)\neq0,\quad \ell\neq0,\quad r\neq0$}\]
\textbf{Supuestos declarados:} r>0; f(r)!=0\par
\section*{Ecuaciones combinadas}
\[\adjustbox{max width=\linewidth}{$\displaystyle E^{\tau}_{\tau}-E^{r}_{r}=0$}\]
\[\adjustbox{max width=\linewidth}{$\displaystyle E^{\tau}_{\tau}-E^{\varphi}_{\varphi}=0$}\]
\[\adjustbox{max width=\linewidth}{$\displaystyle E^{r}_{r}-E^{\varphi}_{\varphi}=0$}\]
\section*{Combinaciones proyectadas y reducidas}
\[\adjustbox{max width=\linewidth}{$\displaystyle E^{\tau}_{\tau}-E^{r}_{r}=0=0$}\]
\[\adjustbox{max width=\linewidth}{$\displaystyle E^{\tau}_{\tau}-E^{\varphi}_{\varphi}=-\frac{\left(-f'\!\left(r\right) + r\,f''\!\left(r\right)\right)}{2\,r}=0$}\]
\[\adjustbox{max width=\linewidth}{$\displaystyle E^{r}_{r}-E^{\varphi}_{\varphi}=-\frac{\left(-f'\!\left(r\right) + r\,f''\!\left(r\right)\right)}{2\,r}=0$}\]
\section*{Ecuaciones absolutas y escalares}
\[\adjustbox{max width=\linewidth}{$\displaystyle E^{\tau}_{\tau}=-\frac{\left(2\,r - {\ell}^{2}\,f'\!\left(r\right)\right)}{2\,{\ell}^{2}\,r}=0$}\]
\[\adjustbox{max width=\linewidth}{$\displaystyle E_\phi=0=0$}\]
Componentes fuera de la diagonal: $0=0$.\par
\section*{Familias verificadas}
Familia 1: verificada en todas las ecuaciones, en el dominio declarado.\par
\[\adjustbox{max width=\linewidth}{$\displaystyle q=0$}\]
\[\adjustbox{max width=\linewidth}{$\displaystyle f\!\left(r\right)=\frac{\left({r}^{2} + C_{1}\,{\ell}^{2}\right)}{{\ell}^{2}}$}\]
\textbf{Método:} ansatz polinomial de grado 2\par
Parámetros libres en esta familia: $\ell,\; C_{1}$.\par
\[\adjustbox{max width=\linewidth}{$\displaystyle \text{Restricciones no nulas:}\qquad C_{1} + \frac{{r}^{2}}{{\ell}^{2}}\neq0,\quad \ell\neq0,\quad r\neq0,\quad \frac{1}{{\ell}^{2}}\neq0$}\]
Supuestos: r>0; f(r)!=0\par
Familia 2: verificada en todas las ecuaciones, en el dominio declarado.\par
\[\adjustbox{max width=\linewidth}{$\displaystyle q=0$}\]
\[\adjustbox{max width=\linewidth}{$\displaystyle f\!\left(r\right)=\frac{\left({r}^{2} + C_{1}\,{\ell}^{2}\right)}{{\ell}^{2}}$}\]
\textbf{Método:} ansatz de potencia con exponente 2\par
Parámetros libres en esta familia: $\ell,\; C_{1}$.\par
\[\adjustbox{max width=\linewidth}{$\displaystyle \text{Restricciones no nulas:}\qquad C_{1} + \frac{{r}^{2}}{{\ell}^{2}}\neq0,\quad \ell\neq0,\quad r\neq0,\quad \frac{1}{{\ell}^{2}}\neq0$}\]
Supuestos: r>0; f(r)!=0\par
Familia 3: verificada en todas las ecuaciones, en el dominio declarado.\par
\[\adjustbox{max width=\linewidth}{$\displaystyle f\!\left(r\right)=\frac{\left({r}^{2} + C_{1}\,{\ell}^{2}\right)}{{\ell}^{2}}$}\]
\textbf{Método:} ansatz polinomial de grado 2\par
Parámetros libres en esta familia: $\ell,\; C_{1},\; q$.\par
\[\adjustbox{max width=\linewidth}{$\displaystyle \text{Restricciones no nulas:}\qquad C_{1} + \frac{{r}^{2}}{{\ell}^{2}}\neq0,\quad \ell\neq0,\quad r\neq0,\quad q\neq0,\quad \frac{1}{{\ell}^{2}}\neq0$}\]
Supuestos: r>0; f(r)!=0\par
Familia 4: verificada en todas las ecuaciones, en el dominio declarado.\par
\[\adjustbox{max width=\linewidth}{$\displaystyle f\!\left(r\right)=\frac{\left({r}^{2} + C_{1}\,{\ell}^{2}\right)}{{\ell}^{2}}$}\]
\textbf{Método:} ansatz de potencia con exponente 2\par
Parámetros libres en esta familia: $\ell,\; C_{1},\; q$.\par
\[\adjustbox{max width=\linewidth}{$\displaystyle \text{Restricciones no nulas:}\qquad C_{1} + \frac{{r}^{2}}{{\ell}^{2}}\neq0,\quad \ell\neq0,\quad r\neq0,\quad q\neq0,\quad \frac{1}{{\ell}^{2}}\neq0$}\]
Supuestos: r>0; f(r)!=0\par
Familia 5: verificada en todas las ecuaciones, en el dominio declarado.\par
\[\adjustbox{max width=\linewidth}{$\displaystyle f\!\left(r\right)=\frac{\left({r}^{2} + C_{1}\,{\ell}^{2}\right)}{{\ell}^{2}}$}\]
\textbf{Método:} Wolfram DSolve (ecuación individual)\par
Parámetros libres en esta familia: $\ell,\; C_{1},\; q$.\par
\[\adjustbox{max width=\linewidth}{$\displaystyle \text{Restricciones no nulas:}\qquad C_{1} + \frac{{r}^{2}}{{\ell}^{2}}\neq0,\quad \ell\neq0,\quad r\neq0$}\]
Supuestos: r>0; f(r)!=0\par
Familia 6: verificada en todas las ecuaciones, en el dominio declarado.\par
\[\adjustbox{max width=\linewidth}{$\displaystyle f\!\left(r\right)=\frac{\left({r}^{2} + C_{1}\,{\ell}^{2}\right)}{{\ell}^{2}}$}\]
\textbf{Método:} Wolfram DSolve (ecuación individual) + Solve de constantes\par
Parámetros libres en esta familia: $\ell,\; C_{1},\; q$.\par
\[\adjustbox{max width=\linewidth}{$\displaystyle \text{Restricciones no nulas:}\qquad C_{1} + \frac{{r}^{2}}{{\ell}^{2}}\neq0,\quad \ell\neq0,\quad r\neq0$}\]
Supuestos: r>0; f(r)!=0\par
\section*{Auditoría de ramas no aceptadas}
Candidatos conservados solo en results.json: rejected=5, undetermined=3, partial=2.\par
Sistema: ODE; orden máximo: 2.\par
Funciones o parámetros libres: q\par
Completitud: no demostrada. Una búsqueda simbólica acotada no certifica que no existan otras ramas diferenciales o singulares.\par
\end{document}